% ********************************************************************
% Ars Classica article
%
% classicthesis's typographic style (Bringhurst-inspired small caps,
% sober rules, Palatino-metric text) run in its "arsclassica" mode —
% Lorenzo Pantieri's re-styling of André Miede's classicthesis, now
% bundled directly into classicthesis.sty as a \usepackage option.
% ********************************************************************
\RequirePackage{silence}
    \WarningFilter{scrartcl}{Usage of package `titlesec'}
    \WarningFilter{titlesec}{Non standard sectioning command detected}
\documentclass[11pt,a4paper]{scrartcl} % KOMA-Script article
\usepackage[T1]{fontenc}
\usepackage[utf8]{inputenc}
\usepackage[style=arsclassica,nochapters]{classicthesis}

\title{\rmfamily\normalfont\spacedallcaps{A Note on Typographic Restraint}}
\author{\spacedlowsmallcaps{Jane Researcher}}
\date{\today}

\begin{document}

\pagestyle{plain}
\maketitle

\begin{abstract}
    \noindent This article demonstrates the \emph{arsclassica} styling of the
    \texttt{classicthesis} package: small-caps headings, a restrained rule
    under the title, and Palatino-metric body text. It is meant as a starting
    point for a short paper or article-length piece of writing, not a full
    thesis.
\end{abstract}

\tableofcontents

\section{Introduction}
\label{sec:introduction}

Classic typography favours restraint: generous margins, a single text
typeface, and small capitals for emphasis rather than boldface. This
template applies that restraint via the \texttt{arsclassica} style bundled
with classicthesis~\cite{classicthesis}, itself inspired by Bringhurst's
\emph{The Elements of Typographic Style}~\cite{bringhurst2004}.

\subsection{Why a KOMA-Script base}
\label{sec:koma}

The document class is \texttt{scrartcl} (KOMA-Script), which classicthesis
was written against; a plain \texttt{article} class also works but loses a
few spacing refinements KOMA provides for headings and margins.

\section{Using this template}
\label{sec:usage}

Replace the title, author and abstract above, then write sections as usual.
Citations use the classic \texttt{bibtex} + \texttt{plain} style pipeline,
already wired to \texttt{references.bib}: for example, this sentence cites
Miede's original thesis template~\cite{miede2012} to show a resolved
reference rather than a \texttt{[?]} mark.

\subsection{Headings}

Section and subsection headings render in small caps via the arsclassica
style. Nest headings up to \verb|\subsubsection|; deeper levels are
unnumbered by classicthesis's own configuration.

\section{Conclusion}
\label{sec:conclusion}

Swap in real content, a real bibliography, and this becomes a short article
or report set in the same style classicthesis uses for full theses.

\bibliographystyle{plain}
\bibliography{references}

\end{document}
